\documentclass[10pt,a4paper]{amsart}
\usepackage{amsmath,amssymb,amsthm,graphicx}
% Retained minimal preamble for the Stacks Project chapter more-morphisms.tex (More on Morphisms).
% Verbatim excerpt of the chapter preamble: Xy-pic setup (native preamble lines 14-18) and the native
% theorem declarations with the mathematical macros (native preamble lines 154-217). Copied mechanically;
% nothing added, removed or redefined. The wrapper class and the amsmath/amssymb/amsthm/graphicx base
% packages are supplied by the assembler, so the class/externaldocument part of the native preamble is omitted.
% For commutative diagrams we use Xy-pic
\usepackage[all]{xy}

% We use 2cell for 2-commutative diagrams.
\xyoption{2cell}
% Theorem environments.
%
\theoremstyle{plain}
\newtheorem{theorem}[subsection]{Theorem}
\newtheorem{proposition}[subsection]{Proposition}
\newtheorem{lemma}[subsection]{Lemma}

\theoremstyle{definition}
\newtheorem{definition}[subsection]{Definition}
\newtheorem{example}[subsection]{Example}
\newtheorem{exercise}[subsection]{Exercise}
\newtheorem{situation}[subsection]{Situation}

\theoremstyle{remark}
\newtheorem{remark}[subsection]{Remark}
\newtheorem{remarks}[subsection]{Remarks}

\numberwithin{equation}{subsection}

% Macros
%
\def\lim{\mathop{\mathrm{lim}}\nolimits}
\def\colim{\mathop{\mathrm{colim}}\nolimits}
\def\Spec{\mathop{\mathrm{Spec}}}
\def\Hom{\mathop{\mathrm{Hom}}\nolimits}
\def\Ext{\mathop{\mathrm{Ext}}\nolimits}
\def\SheafHom{\mathop{\mathcal{H}\!\mathit{om}}\nolimits}
\def\SheafExt{\mathop{\mathcal{E}\!\mathit{xt}}\nolimits}
\def\Sch{\mathit{Sch}}
\def\Mor{\mathop{\mathrm{Mor}}\nolimits}
\def\Ob{\mathop{\mathrm{Ob}}\nolimits}
\def\Sh{\mathop{\mathit{Sh}}\nolimits}
\def\NL{\mathop{N\!L}\nolimits}
\def\CH{\mathop{\mathrm{CH}}\nolimits}
\def\proetale{{pro\text{-}\acute{e}tale}}
\def\etale{{\acute{e}tale}}
\def\QCoh{\mathit{QCoh}}
\def\Ker{\mathop{\mathrm{Ker}}}
\def\Im{\mathop{\mathrm{Im}}}
\def\Coker{\mathop{\mathrm{Coker}}}
\def\Coim{\mathop{\mathrm{Coim}}}

% Boxtimes
%
\DeclareMathSymbol{\boxtimes}{\mathbin}{AMSa}{"02}

%
% Macros for moduli stacks/spaces
%
\def\QCohstack{\mathcal{QC}\!\mathit{oh}}
\def\Cohstack{\mathcal{C}\!\mathit{oh}}
\def\Spacesstack{\mathcal{S}\!\mathit{paces}}
\def\Quotfunctor{\mathrm{Quot}}
\def\Hilbfunctor{\mathrm{Hilb}}
\def\Curvesstack{\mathcal{C}\!\mathit{urves}}
\def\Polarizedstack{\mathcal{P}\!\mathit{olarized}}
\def\Complexesstack{\mathcal{C}\!\mathit{omplexes}}
% \Pic is the operator that assigns to X its picard group, usage \Pic(X)
% \Picardstack_{X/B} denotes the Picard stack of X over B
% \Picardfunctor_{X/B} denotes the Picard functor of X over B
\def\Pic{\mathop{\mathrm{Pic}}\nolimits}
\def\Picardstack{\mathcal{P}\!\mathit{ic}}
\def\Picardfunctor{\mathrm{Pic}}
\def\Deformationcategory{\mathcal{D}\!\mathit{ef}}
\newenvironment{reference}{}{}
\newenvironment{slogan}{}{}
\newenvironment{history}{}{}
\usepackage{cleveref}
\crefname{originallemma}{original lemma}{original lemmas}
\Crefname{originallemma}{Original lemma}{Original lemmas}
\makeatletter
\newlabel{algebra-lemma-base-change-flat-up-down}{{algebra-lemma-base-change-flat-up-down}{source}{}{sourcepointer-1}{}}
\newlabel{algebra-lemma-base-change-flat-up-down@cref}{{[originallemma][0][]algebra-lemma-base-change-flat-up-down}{[1][0][]source}{}{}{}}
\makeatother
\crefname{originallemma}{original lemma}{original lemmas}
\Crefname{originallemma}{Original lemma}{Original lemmas}
\makeatletter
\newlabel{varieties-lemma-CM-base-change}{{varieties-lemma-CM-base-change}{source}{}{sourcepointer-2}{}}
\newlabel{varieties-lemma-CM-base-change@cref}{{[originallemma][0][]varieties-lemma-CM-base-change}{[1][0][]source}{}{}{}}
\makeatother
\crefname{originallemma}{original lemma}{original lemmas}
\Crefname{originallemma}{Original lemma}{Original lemmas}
\makeatletter
\newlabel{varieties-lemma-locally-Noetherian-base-change}{{varieties-lemma-locally-Noetherian-base-change}{source}{}{sourcepointer-3}{}}
\newlabel{varieties-lemma-locally-Noetherian-base-change@cref}{{[originallemma][0][]varieties-lemma-locally-Noetherian-base-change}{[1][0][]source}{}{}{}}
\makeatother
\title{Cohen-Macaulayness of morphisms of schemes is stable under flat base change}
\author{The Stacks Project Authors}
\date{stacks-project:more-morphisms.tex:local-snapshot-20261011; GNU FDL 1.2 or later}
\begin{document}
\maketitle
\noindent\textit{Editorial note.} Statement span is lines 5911-5926 of sources/stacks-more-morphisms/pinned-original-source.tex; proof span is lines 5928-5978, the complete author proof (49 body lines).\par\medskip
\noindent\textit{Editorial note.} Navigation stubs: 3 cited-but-not-included labels are supplied as non-load-bearing pointers with their pinned tags where the snapshot tag index has them.\par\medskip
\noindent\textit{Editorial note.} extra\_packages is empty on purpose: the retained preamble already loads xy with options and a second declaration clashes in the assembler wrapper.\par\medskip
\noindent\textit{Editorial note.} Text is retained verbatim, including the author's own typos and placeholders.\par\medskip
\noindent\textit{Allowed standard prerequisites (editorial).}\begin{itemize}
\item Explicit prerequisite boundary. Cohen-Macaulay morphisms, depth and flat local ring maps, and base change of morphisms; the transfer is retained.
\end{itemize}\medskip
\ifdefined\hypertarget\hypertarget{sourcepointer-1}{}\fi
\noindent\textit{Source reference.} Original lemma algebra-lemma-base-change-flat-up-down: sources/stacks-more-morphisms/pinned-original-source.tex, source lines 5978--5978. Cited by the retained statement or proof but not part of the collected spans; supplied only as a navigation pointer. The argument does not depend on its content.\par\medskip
\ifdefined\hypertarget\hypertarget{sourcepointer-2}{}\fi
\noindent\textit{Source reference.} Original lemma varieties-lemma-CM-base-change: sources/stacks-more-morphisms/pinned-original-source.tex, source lines 5978--5978. Cited by the retained statement or proof but not part of the collected spans; supplied only as a navigation pointer. The argument does not depend on its content.\par\medskip
\ifdefined\hypertarget\hypertarget{sourcepointer-3}{}\fi
\noindent\textit{Source reference.} Original lemma varieties-lemma-locally-Noetherian-base-change: sources/stacks-more-morphisms/pinned-original-source.tex, source lines 5978--5978. Cited by the retained statement or proof but not part of the collected spans; supplied only as a navigation pointer. The argument does not depend on its content.\par\medskip
\begin{lemma}
\label{lemma-base-change-CM}
Let $f : X \to Y$ be a morphism of schemes.
Assume that all the fibres $X_y$ are locally Noetherian schemes.
Let $Y' \to Y$ be locally of finite type. Let $f' : X' \to Y'$
be the base change of $f$.
Let $x' \in X'$ be a point with image $x \in X$.
\begin{enumerate}
\item If $f$ is Cohen-Macaulay at $x$, then
$f' : X' \to Y'$ is Cohen-Macaulay at $x'$.
\item If $f$ is flat at $x$ and $f'$ is Cohen-Macaulay at $x'$, then $f$
is Cohen-Macaulay at $x$.
\item If $Y' \to Y$ is flat at $f'(x')$ and $f'$ is Cohen-Macaulay at
$x'$, then $f$ is Cohen-Macaulay at $x$.
\end{enumerate}
\end{lemma}

\begin{proof}
Note that the assumption on $Y' \to Y$ implies that for $y' \in Y'$
mapping to $y \in Y$ the field extension $\kappa(y')/\kappa(y)$
is finitely generated. Hence also all the fibres
$X'_{y'} = (X_y)_{\kappa(y')}$ are locally Noetherian, see
Varieties, Lemma \ref{varieties-lemma-locally-Noetherian-base-change}.
Thus the lemma makes sense. Set $y' = f'(x')$ and $y = f(x)$.
Hence we get the following commutative diagram of local rings
$$
\xymatrix{
\mathcal{O}_{X', x'} & \mathcal{O}_{X, x} \ar[l] \\
\mathcal{O}_{Y', y'} \ar[u] & \mathcal{O}_{Y, y} \ar[l] \ar[u]
}
$$
where the upper left corner is a localization of the tensor product
of the upper right and lower left corners over the lower right corner.

\medskip\noindent
Assume $f$ is Cohen-Macaulay at $x$.
The flatness of $\mathcal{O}_{Y, y} \to \mathcal{O}_{X, x}$
implies the flatness of $\mathcal{O}_{Y', y'} \to \mathcal{O}_{X', x'}$, see
Algebra, Lemma \ref{algebra-lemma-base-change-flat-up-down}.
The fact that $\mathcal{O}_{X, x}/\mathfrak m_y\mathcal{O}_{X, x}$
is Cohen-Macaulay implies that
$\mathcal{O}_{X', x'}/\mathfrak m_{y'}\mathcal{O}_{X', x'}$
is Cohen-Macaulay, see
Varieties, Lemma \ref{varieties-lemma-CM-base-change}. Hence we see that $f'$
is Cohen-Macaulay at $x'$.

\medskip\noindent
Assume $f$ is flat at $x$ and $f'$ is Cohen-Macaulay at $x'$.
The fact that $\mathcal{O}_{X', x'}/\mathfrak m_{y'}\mathcal{O}_{X', x'}$
is Cohen-Macaulay implies that
$\mathcal{O}_{X, x}/\mathfrak m_y\mathcal{O}_{X, x}$
is Cohen-Macaulay, see
Varieties, Lemma \ref{varieties-lemma-CM-base-change}.
Hence we see that $f$ is Cohen-Macaulay at $x$.

\medskip\noindent
Assume $Y' \to Y$ is flat at $y'$ and $f'$ is Cohen-Macaulay at
$x'$. The flatness of $\mathcal{O}_{Y', y'} \to \mathcal{O}_{X', x'}$
and $\mathcal{O}_{Y, y} \to \mathcal{O}_{Y', y'}$ implies the flatness
of $\mathcal{O}_{Y, y} \to \mathcal{O}_{X, x}$, see
Algebra, Lemma \ref{algebra-lemma-base-change-flat-up-down}.
The fact that $\mathcal{O}_{X', x'}/\mathfrak m_{y'}\mathcal{O}_{X', x'}$
is Cohen-Macaulay implies that
$\mathcal{O}_{X, x}/\mathfrak m_y\mathcal{O}_{X, x}$
is Cohen-Macaulay, see
Varieties, Lemma \ref{varieties-lemma-CM-base-change}. Hence we see that $f$
is Cohen-Macaulay at $x$.
\end{proof}

\end{document}
